\subsection{Testing the marginal constraints}
\label{subsec:optimal-marginal-verification}

Determination is an exact statement.  To use the marginals as a certificate,
we also need to know how reliably they detect a deviation from the target.
Sampling the \(m\) projectors above uniformly gives a worst-case rejection
rate of only \(1/m\) on states orthogonal to \(\psi\).  We can improve
this by using the complementary half as well.

Fix complementary halves \(B,C\), and put
\[
 \mathcal S=\{B\cup\{j\}:j\in C\}
       \cup\{C\cup\{i\}:i\in B\},\qquad
 \Pi_S=q^{m-1}\rho_S^\psi\otimes I_{S^c}.
\]
In one round, choose \(S\) uniformly from these \(2m\) supports and
measure the binary test \(\{\Pi_S,I-\Pi_S\}\), accepting on \(\Pi_S\).
This requires access to \(m+1\) parties of one copy.  Joint operations
within that subset are allowed.  The restriction is on the \emph{total
support} of a test; it differs from measuring all parties separately and
combining their outcomes.  A binary test need not be one laboratory
measurement basis.

The verification-operator and spectral-gap framework is standard
\cite{PallisterLindenMontanaro2018,DangniamHanZhu2020}, as are fidelity
certificates from stabilizer generators and frustration-free Hamiltonians
\cite{KalevKyrillidisLinke2019,ZhuLiChen2024}.  Setting minimization and
qudit stabilizer witnesses have also been studied
\cite{LiZhangLiZhu2021,SzczepaniakMakutaAugusiak2026}.
Here the AME support geometry lets us calculate the optimal rejection
rate under the stated support restriction.

\begin{proposition}[Complementary marginal tests]
\label{prop:complementary-marginal-tests}
Let \(\psi\) be a stabilizer \(\AME(2m,q)\) state with \(m\geq2\).
The verification operator and its gap are
\begin{equation}
 \Omega_* = \frac1{2m}\sum_{S\in\mathcal S}\Pi_S,
 \qquad \nu_* = \frac{m+1}{2m},\qquad
 I-\Omega_*\geq\nu_*(I-|\psi\rangle\langle\psi|).
 \label{eq:complementary-gap}
\end{equation}
The coefficient \(\nu_*\) is optimal among randomized perfectly complete
binary tests, each supported on at most \(m+1\) parties.  For every density
operator \(\sigma\), its rejection probability
\(r(\sigma)=\Tr[(I-\Omega_*)\sigma]\) satisfies
\begin{equation}
 1-\langle\psi|\sigma|\psi\rangle
 \leq\frac{r(\sigma)}{\nu_*}
 \leq\frac1{m+1}\sum_{S\in\mathcal S}
       \frac12\|\sigma_S-\rho_S^\psi\|_1.
 \label{eq:complementary-fidelity}
\end{equation}
\end{proposition}

\begin{proof}
The key is that \emph{any} \(m\) of the supported spaces \(L(S)\)
form a direct sum equal to \(L\).  To see this, select
\(I\subset B,J\subset C\) with \(|I|+|J|=m\).  In a relation
\(\sum_{i\in I}u_i=-\sum_{j\in J}v_j\), with
\(u_i\in L(C\cup\{i\})\) and \(v_j\in L(B\cup\{j\})\), the common
sum is supported in \(I\cup J\).  The support theorem makes it zero.
Projection to \(i\) isolates \(u_i\) and is injective on its subgroup,
so every term vanishes; likewise for \(v_j\).  Each subgroup has dimension
\(2e\) over \(\F_p\), so the dimensions add to \(\dim L\).

Use the actual stabilizer lifts fixing \(\psi\).  Their joint eigenbasis
is indexed by characters of \(L\).  A projector \(\Pi_S\) accepts
exactly those characters trivial on \(L(S)\).  A nontrivial character
cannot pass \(m\) tests, hence fails at least \(m+1\).  This proves the
operator inequality.  It is sharp: any \(m-1\) chosen subgroups have a
nonzero annihilating character, which passes exactly those tests.

For optimality in the larger class of allowed tests, a perfectly complete
effect \(E\) satisfies \(E\psi=\psi\).  A traceless one-site unitary
\(V_i\) makes \(V_i\psi\) orthogonal to \(\psi\), and every test omitting
\(i\) still accepts it.  Thus the gap is at most the probability of
including party \(i\).  Averaging over \(2m\) parties bounds that
probability, for some \(i\), by \((m+1)/(2m)\).
Finally take the expectation of the operator inequality.  Testing a
marginal discrepancy against its support projector bounds each rejection
probability by \(\tfrac12\|\sigma_S-\rho_S^\psi\|_1\); average and divide
by \(\nu_*\).
\end{proof}

For the four-qutrit state of Section~\ref{subsec:four-qutrit-example},
the two tests on \(123\) and \(124\) have gap \(1/2\); using all four
three-party tests raises it to \(3/4\).  The state
\(Z_1(1)\psi\) passes the test omitting party 1 and fails the other
three: on each support containing party 1, surjectivity onto its Weyl
labels makes the induced character nontrivial.  These tests detect exactly
the character change that left the transition maps unchanged in the example.

Only acceptance probabilities are needed for the first fidelity bound;
full marginal tomography is unnecessary.  If the prepared state is promised
to be \(U\psi\) for a product unitary \(U\), the same certificate enters
the robust rigidity theorem through
Corollary~\ref{cor:marginal-certified-rounding}.
Appendix~\ref{app:marginal-verification} gives the exact tradeoff for fewer
or unequally weighted tests and extends the uniform gap to every AME state.
